% ============================================================
% main_v4_additions.tex
% Copy-ready additions for main_v3.tex
%
% This file is not intended to compile by itself.
% Copy each marked block into the indicated location in main_v3.tex.
% ============================================================


% ============================================================
% BLOCK A
% INSERT AT THE END OF:
% \subsection{Overview of the transformation}
% ============================================================

\paragraph{Two nonredundant congestion branches.}
The demand--capacity ratio does not determine emissions through queue speed
alone. Instead, it generates two distinct traffic-state branches:
\begin{equation}
\label{eq:two_branch_state}
z=\frac{D}{C}
\quad\longrightarrow\quad
\begin{cases}
P(z),\,w_{t_2}(z),\,\bar w(z),
& \text{delay and exposure branch},\\[1mm]
\mu(z)\longrightarrow v_q(z)\longrightarrow\Gamma(z),
& \text{physical state and emission-severity branch}.
\end{cases}
\end{equation}
The first branch determines how much excess-time exposure vehicles experience.
The second determines the emission consequence per unit of that exposure.
The two branches meet in
\begin{equation}
\label{eq:two_branch_emission_summary}
e(z)
=
r(v_{\mathrm{ref}})T_c
+
\Gamma(z)\bar w(z),
\qquad
\Gamma(z)
\equiv
\Gamma\!\left(v_q(z)\right).
\end{equation}
Thus, the proposed model is neither a delay-only model nor a speed-only model.
It is a joint exposure--severity model.


% ============================================================
% BLOCK B
% REPLACE THE CURRENT PIPELINE FIGURE OR INSERT AFTER IT.
% This is a compile-safe placeholder for the revised Figure 1.
% ============================================================

\begin{figure}[t]
\centering
\fbox{%
\begin{minipage}{0.94\textwidth}
\centering
\vspace{3mm}
\textbf{PLACEHOLDER FOR REVISED FIGURE 1}

\vspace{3mm}
\[
z=\frac{D}{C}
\]

\[
\begin{array}{ccccc}
& \nearrow &
P(z),\,w_{t_2}(z),\,\bar w(z)
& \quad &
\text{delay/exposure branch}
\\[2mm]
z
&&&&
\\[-1mm]
& \searrow &
\mu(z)\rightarrow v_q(z)\rightarrow\Gamma(z)
&&
\text{physical/emission-severity branch}
\end{array}
\]

\[
\left[\bar w(z),\Gamma(z)\right]
\longrightarrow
e(z)
=
r(v_{\mathrm{ref}})T_c+\Gamma(z)\bar w(z).
\]

\[
\left[w(\tau;z),v_q(z)\right]
\longrightarrow
t_Q(\tau;z)
\longrightarrow
EC_1(\tau;z).
\]

\vspace{3mm}
\end{minipage}
}
\caption{Two-branch structure of the proposed queue-emission transformation.
The delay branch determines the quantity of congestion exposure, whereas the
discharge--FD branch determines the emission consequence per unit of exposure.
For a fixed episode conditioned on $z$, $\mu(z)$, $v_q(z)$, and $\Gamma(z)$
are constant, while $w(\tau;z)$ varies among vehicles.}
\label{fig:two_branch_placeholder}
\end{figure}


% ============================================================
% BLOCK C
% INSERT AT THE BEGINNING OF:
% \section{Two-regime queue traversal and emission foundation}
% BEFORE:
% \subsection{Lawson's distinction between delay and time in queue}
% ============================================================

\subsection{Classical fixed-bottleneck benchmark}
\label{subsec:fixed_bottleneck}

The proposed construction is first interpreted relative to a classical
fixed bottleneck. Let $\mu_0$ denote the constant discharge rate of an active
bottleneck, and let $\lambda(\tau)$ denote the arrival rate to the queue.
While the queue exists, its accumulation satisfies
\begin{equation}
\label{eq:fixed_bottleneck_accumulation}
\dot q(\tau)=\lambda(\tau)-\mu_0.
\end{equation}
For the elementary two-rate representation,
\begin{equation}
\label{eq:two_arrival_rates}
\lambda(\tau)=
\begin{cases}
\lambda_H>\mu_0, & t_0\leq\tau<t_2,\\[1mm]
\lambda_L<\mu_0, & t_2<\tau\leq t_3.
\end{cases}
\end{equation}
The high-arrival branch accumulates the queue, whereas the low-arrival branch
dissipates it. For a vertical point queue, the virtual waiting time is
$w(\tau)=q(\tau)/\mu_0$, so
\begin{equation}
\label{eq:waiting_slopes_fixed_bottleneck}
\frac{dw(\tau)}{d\tau}
=
\begin{cases}
(\lambda_H-\mu_0)/\mu_0>0,
& t_0\leq\tau<t_2,\\[2mm]
(\lambda_L-\mu_0)/\mu_0<0,
& t_2<\tau\leq t_3.
\end{cases}
\end{equation}

This benchmark separates the temporal demand-side mechanism from the
bottleneck supply state. The two arrival-rate branches determine how waiting
time accumulates and dissipates, while the bottleneck discharge $\mu_0$
remains fixed. Consequently, an FD-equivalent congested vehicle speed
reconstructed from $\mu_0$ is also fixed within the representative episode.

\paragraph{Reduced-form schedule-early and schedule-late interpretation.}
The present paper does not observe or optimize a complete set of link-level
departure times. It therefore does not claim to reconstruct a full endogenous
Vickrey departure-time equilibrium. Nevertheless, the aggregate queue episode
admits a corner-point scheduling interpretation.

Let $A$ denote the unit disutility of waiting time, $B$ the unit schedule-early
penalty, and $C_s$ the unit schedule-late penalty. The generalized scheduling
disutility is
\begin{equation}
\label{eq:ABC_schedule_cost}
U(\tau)
=
A w(\tau)
+
B S_E(\tau)
+
C_s S_L(\tau).
\end{equation}
Under the symmetric reduced-form case used by the PAQ profile,
\begin{equation}
t_2-t_0=t_3-t_2=\frac{P}{2},
\qquad
B=C_s.
\end{equation}
The two zero-wait endpoints and the maximum-waiting point then admit the
corner-point consistency condition
\begin{equation}
\label{eq:corner_point_consistency}
B\frac{P}{2}
=
A w_{t_2}
=
C_s\frac{P}{2}.
\end{equation}
The beginning-of-episode traveler experiences no queueing but incurs
schedule-early displacement; the traveler associated with $t_2$ experiences
the maximum waiting time but no schedule displacement; and the
end-of-episode traveler experiences no queueing but incurs schedule-late
displacement.

Equation~\eqref{eq:corner_point_consistency} is used only as a reduced-form
behavioral interpretation of the observed congestion duration and maximum
waiting time. Establishing equal generalized cost at every used departure
time would require a full dynamic departure-time model, which is outside the
scope of the present paper.

\paragraph{Behavioral interpretation of the PAQ scaling relation.}
Using
\begin{equation}
w_{t_2}(z)=T_c f_p[P(z)]^s,
\end{equation}
the symmetric corner-point relation implies
\begin{equation}
\label{eq:ABC_qvdf_interpretation}
\frac{B}{A}
=
\frac{C_s}{A}
=
\frac{2w_{t_2}(z)}{P(z)}
=
2T_c f_p[P(z)]^{s-1}.
\end{equation}
When $s=1$, the implied waiting-versus-schedule trade-off is constant across
episodes. When $s\neq1$, $f_p$ and $s$ should be interpreted as reduced-form
queue-shape parameters rather than fixed traveler-preference coefficients.
They may jointly reflect scheduling behavior, heterogeneity, temporal
aggregation, and facility-specific operating conditions.


% ============================================================
% BLOCK D
% OPTIONAL FIGURE PLACEHOLDER.
% INSERT AFTER BLOCK C OR AFTER THE EXISTING VEHICLE TRAJECTORY FIGURE.
% ============================================================

\begin{figure}[t]
\centering
\fbox{%
\begin{minipage}{0.94\textwidth}
\centering
\vspace{3mm}
\textbf{PLACEHOLDER FOR CLASSICAL BOTTLENECK FIGURE}

\vspace{2mm}
\textbf{Panel A:} cumulative arrivals and departures, with
$\lambda_H>\mu_0$ during accumulation,
$\lambda_L<\mu_0$ during dissipation, and constant departure slope $\mu_0$.

\vspace{2mm}
\textbf{Panel B:} waiting/schedule trade-off:
\[
w(t_0)=0,\qquad
w(t_2)=w_{t_2},\qquad
w(t_3)=0,
\]
and
\[
B\frac{P}{2}=Aw_{t_2}=C_s\frac{P}{2}.
\]

\vspace{2mm}
\textbf{Panel C:} incremental individual emission:
\[
\Delta e(\tau)=\Gamma_0 w(\tau),
\]
shown as a vertical rescaling of the waiting-time profile.

\vspace{3mm}
\end{minipage}
}
\caption{Classical fixed-bottleneck benchmark and its emission augmentation.
The arrival rate changes across the accumulation and dissipation branches,
whereas the bottleneck discharge, FD-equivalent queued vehicle speed, and
delay-to-emission factor remain constant within the episode.}
\label{fig:fixed_bottleneck_placeholder}
\end{figure}


% ============================================================
% BLOCK E
% INSERT IMMEDIATELY AFTER EQUATION:
% \eqref{eq:single_vehicle_emission}
% ============================================================

\paragraph{Interpretation of the delay-to-emission factor.}
The factor $\Gamma$ is a per-vehicle, per-unit-delay emission conversion
coefficient. It is not a queue-accumulation variable and is not intrinsically
time dependent. For a fixed bottleneck discharge $\mu_0$, define
\begin{equation}
\label{eq:fixed_mu_vq_gamma}
v_{q,0}
=
V_{\mathrm{FD}}^{\mathrm{cong}}(\mu_0),
\qquad
\Gamma_0
=
\Gamma(v_{q,0}).
\end{equation}
Because $\mu_0$ is constant within the episode, $v_{q,0}$ and $\Gamma_0$ are
also constant. The temporal variation in individual congestion emissions is
generated by the vehicle-specific delay:
\begin{equation}
\label{eq:fixed_bottleneck_vehicle_emission}
EC_1(\tau)
=
r(v_{\mathrm{ref}})T_c
+
\Gamma_0 w(\tau).
\end{equation}
Hence, under the classical fixed-discharge benchmark, the incremental
emission profile is an exact vertical rescaling of the waiting-time profile:
\begin{equation}
\label{eq:fixed_bottleneck_incremental_emission}
\Delta e(\tau)
\equiv
EC_1(\tau)-r(v_{\mathrm{ref}})T_c
=
\Gamma_0 w(\tau).
\end{equation}

In the QVDF generalization, $\Gamma$ becomes a function of the congestion
index only when different representative episodes are compared:
\begin{equation}
\label{eq:gamma_as_z_function}
\Gamma(z)
=
\Gamma\!\left(
V_{\mathrm{FD}}^{\mathrm{cong}}[\mu(z)]
\right).
\end{equation}
Within any one episode conditioned on $z$, $\mu(z)$, $v_q(z)$, and
$\Gamma(z)$ remain constant under the analytical reconstruction, while
$w(\tau;z)$ varies among vehicle entry times. Thus, $\Gamma(z)$ controls
emission severity per unit of delay, whereas $w(\tau;z)$ controls the
amount of vehicle-specific exposure.


% ============================================================
% BLOCK F
% INSERT IN SECTION:
% "From individual traversals to a congested episode"
% AFTER THE DEFINITION OF ARRIVAL-WEIGHTED MEAN DELAY.
% ============================================================

\paragraph{Cumulative-curve interpretation under fixed discharge.}
For the classical fixed-discharge benchmark, $\Gamma_0$ is constant and may
be taken outside the episode integral. Consequently, the incremental
bottleneck emission is
\begin{equation}
\label{eq:fixed_bottleneck_aggregate_emission}
EC_{\mathrm{inc}}
=
\Gamma_0
\int_{\mathcal T}
a(\tau)w(\tau)\,d\tau
=
\Gamma_0 D\bar w.
\end{equation}
The integral in Equation~\eqref{eq:fixed_bottleneck_aggregate_emission} is the
total vehicle excess travel time represented by the area between the virtual
arrival and departure curves in the cumulative input--output diagram.
Therefore, the proposed physicalization does not replace the classical
cumulative curves; it assigns an environmental consequence to their
vehicle-time area.


% ============================================================
% BLOCK G
% INSERT AT THE BEGINNING OF:
% \section{Traffic-state reconstruction from the FD--QVDF}
% AFTER THE OPENING PARAGRAPH.
% ============================================================

\subsection{Nested interpretation of the discharge model}
\label{subsec:nested_discharge}

The classical benchmark treats the bottleneck discharge $\mu_0$ as constant
within and across representative congestion episodes. The QVDF extension
retains constancy within each episode but permits the episode-equivalent
discharge to vary across episodes indexed by $z$. Accordingly,
$\mu(z)$ should not be interpreted as an instantaneously varying service rate
inside one queue. Rather, it is the constant effective discharge assigned to
a representative episode with congestion level $z$.

This interpretation produces three nested cases:
\begin{enumerate}[label=(\alph*),leftmargin=2em]
\item \textbf{Classical fixed bottleneck:}
$\mu=\mu_0$, so $v_q=v_{q,0}$ and $\Gamma=\Gamma_0$.
\item \textbf{Episode-specific bottleneck:}
each episode $j$ has a constant $\mu_j$, $v_{q,j}$, and $\Gamma_j$, but those
values may differ across episodes.
\item \textbf{QVDF-indexed bottleneck:}
the episode-specific values are represented by
$\mu_j\approx\mu(z_j)$, giving
$v_q(z)=V_{\mathrm{FD}}^{\mathrm{cong}}[\mu(z)]$ and
$\Gamma(z)=\Gamma(v_q(z))$.
\end{enumerate}
The first case is the benchmark, the second is the empirical intermediate
representation, and the third is the planning-scale closure developed in
this paper.


% ============================================================
% BLOCK H
% INSERT AFTER EQUATION:
% \eqref{eq:qvdf_duration_discharge}
% ============================================================

\paragraph{Within-episode constancy versus cross-episode variation.}
Equation~\eqref{eq:qvdf_duration_discharge} indexes an
episode-equivalent discharge by the aggregate congestion state. It does not
imply that the discharge changes continuously with time during a single
episode. Conditional on a given $z$, the analytical reconstruction assumes
\begin{equation}
\mu(\tau\mid z)=\mu(z),
\qquad
t_0\leq\tau\leq t_3.
\end{equation}
The temporal queue profile within that episode is generated by the difference
between the time-dependent arrival rate and the fixed episode discharge.
Across representative episodes, however,
\begin{equation}
z_i\neq z_j
\quad\Longrightarrow\quad
\mu(z_i)\ \text{may differ from}\ \mu(z_j).
\end{equation}
Thus, the proposed $\mu(z)$ relation is a cross-episode QVDF hypothesis,
with the classical constant-discharge model nested as a special case.


% ============================================================
% BLOCK I
% INSERT AFTER EQUATION:
% \eqref{eq:vq_fd_qvdf}
% ============================================================

\paragraph{Meaning of the reconstructed speed.}
The quantity $v_q(z)$ is an episode-equivalent vehicle operating speed on the
congested branch of the fundamental diagram. It is not the propagation speed
of the upstream or downstream queue boundary, a kinematic shockwave speed,
the speed of a moving bottleneck, or the speed at which the spatial queue
length changes. The model does not reconstruct the full spatial evolution of
the queue. Instead, it supplies the within-queue vehicle speed needed to
evaluate the queued-regime emission rate in the two-regime trajectory model.


% ============================================================
% BLOCK J
% REPLACE OR EXPAND THE PARAGRAPH DISCUSSING n=1 AND n>1
% IN "Admissibility of the discharge--speed chain".
% ============================================================

The boundary case $n=1$ yields
\begin{equation}
\mu(z)=\frac{C}{f_d},
\end{equation}
which is constant across representative congestion episodes. It therefore
recovers the classical fixed-discharge benchmark, subject to the adopted
normalization. When $n>1$, the model permits the episode-equivalent discharge
to decline across episodes as the aggregate oversaturation index increases.
This is a cross-episode generalization, not a claim that the bottleneck
service rate necessarily declines continuously during one congestion episode.
A calibration with $n<1$ would imply that the reconstructed episode discharge
increases with congestion and is therefore excluded from the physical
discharge--speed interpretation used here.


% ============================================================
% BLOCK K
% INSERT IN:
% \section{Assignment compatibility}
% BEFORE:
% \subsection{Generalized link cost}
% ============================================================

\subsection{Three distinct uses of the emission function}
\label{subsec:three_uses}

The same queue-emission function supports three conceptually distinct uses.

\paragraph{Ex-post environmental accounting.}
The observed or assigned traffic pattern is held fixed, and total emissions
are evaluated using
\begin{equation}
EC(z)=D e(z).
\end{equation}
In this use, emissions do not affect route or departure-time choices.

\paragraph{Emission-weighted static assignment.}
Per-vehicle emissions enter a separable link cost together with travel time.
This is the assignment use developed in the present section.

\paragraph{Departure-time DTA extension.}
A future dynamic model could augment the scheduling disutility by an
emission term:
\begin{equation}
\label{eq:emission_augmented_departure_cost}
U_e(\tau;z)
=
A w(\tau;z)
+
B S_E(\tau)
+
C_s S_L(\tau)
+
\eta_e\Gamma(z)w(\tau;z).
\end{equation}
For a fixed episode conditioned on $z$, this becomes
\begin{equation}
\label{eq:effective_waiting_coefficient}
U_e(\tau;z)
=
A_{\mathrm{eff}}(z)w(\tau;z)
+
B S_E(\tau)
+
C_s S_L(\tau),
\qquad
A_{\mathrm{eff}}(z)
=
A+\eta_e\Gamma(z).
\end{equation}
Thus, under a fixed episode state, emissions modify the effective coefficient
on waiting time rather than introducing a new temporal shape. The present
paper does not solve the endogenous departure-time equilibrium implied by
Equation~\eqref{eq:emission_augmented_departure_cost}; it supplies the
traffic-state and emission terms required by such an extension.


% ============================================================
% BLOCK L
% REPLACE THE CURRENT:
% \subsection{Monotonicity condition}
% WITH THIS EXPANDED VERSION.
% ============================================================

\subsection{Monotonicity and convex assignment potential}

Define
\begin{equation}
\Gamma_z(z)
\equiv
\Gamma(v_q(z)).
\end{equation}
Using
\begin{equation}
e(z)
=
r(v_{\mathrm{ref}})T_c
+
\Gamma_z(z)\bar w(z),
\end{equation}
the derivative of the complete generalized link cost is
\begin{equation}
\label{eq:exact_generalized_cost_derivative}
g'(z)
=
\left[
\theta_t+\theta_e\Gamma_z(z)
\right]\bar w'(z)
+
\theta_e\Gamma_z'(z)\bar w(z).
\end{equation}
The exact assignment condition is therefore
\begin{equation}
\label{eq:exact_monotonicity}
g'(z)\geq0
\qquad
\text{on the calibrated assignment domain}.
\end{equation}
The factor $\Gamma_z(z)$ need not itself be linear or convex. What matters
for a convex separable assignment formulation is monotonicity of the complete
generalized link cost.

Under the admissible discharge--speed chain,
$v_q'(z)\leq0$, and
\begin{equation}
\Gamma_z'(z)
=
\Gamma'(v_q(z))v_q'(z).
\end{equation}
A convenient sufficient condition for the emission component to be
nondecreasing is
\begin{equation}
\label{eq:gamma_sufficient_condition_v4}
\Gamma(v_q)\geq0,
\qquad
\Gamma'(v_q)\leq0
\end{equation}
for all queued speeds in the assignment domain. This condition is sufficient
but not necessary: the positive travel-time term in
Equation~\eqref{eq:exact_generalized_cost_derivative} may preserve
$g'(z)\geq0$ even when the emission component alone is not monotone.

Let $x$ denote link demand and let $K$ denote the corresponding period
capacity, so that
\begin{equation}
z=\frac{x}{K}.
\end{equation}
If $g(z)$ is continuous and nondecreasing, the link contribution to the
Beckmann potential is
\begin{equation}
\label{eq:flow_space_beckmann}
\Phi(x)
=
\int_0^x
g\!\left(\frac{u}{K}\right)\,du
=
K\int_0^{x/K}g(s)\,ds.
\end{equation}
Where derivatives exist,
\begin{equation}
\label{eq:beckmann_convexity_proof}
\Phi''(x)
=
\frac{1}{K}
g'\!\left(\frac{x}{K}\right)
\geq0.
\end{equation}
Hence, convexity of the Beckmann potential follows from monotonicity of the
complete generalized cost. It does not require $\Gamma(z)$ itself to be
convex.


% ============================================================
% BLOCK M
% INSERT BEFORE THE CURRENT POLYNOMIAL PRIMITIVE EQUATION.
% ============================================================

An integral with respect to $z$ is a normalized primitive rather than the
full flow-space link potential. Using $x=Kz$, the corresponding Beckmann
contribution is
\begin{equation}
\label{eq:normalized_to_flow_primitive}
\int_{x_a}^{x_b}g\!\left(\frac{u}{K}\right)\,du
=
K\int_{z_a}^{z_b}g(z)\,dz.
\end{equation}
If
\begin{equation}
\Gamma_z(z)
\approx
F(z)
=
\sum_k a_k z^k,
\end{equation}
then the normalized congestion-emission primitive remains
\begin{equation}
\int_{z_a}^{z_b}
F(z)\bar w(z)\,dz
=
T_c\alpha
\sum_k
\frac{a_k}{k+\beta+1}
\left[
z^{k+\beta+1}
\right]_{z_a}^{z_b},
\end{equation}
and the corresponding flow-space contribution is obtained by multiplying
this expression by $K$.


% ============================================================
% BLOCK N
% OPTIONAL FIGURE PLACEHOLDER FOR FUTURE ATLANTA VALIDATION.
% INSERT IN THE NUMERICAL SECTION WHERE APPROPRIATE.
% ============================================================

\begin{figure}[t]
\centering
\fbox{%
\begin{minipage}{0.94\textwidth}
\centering
\vspace{3mm}
\textbf{PLACEHOLDER FOR EXTERNAL DISCHARGE VALIDATION}

\vspace{2mm}
Suggested contents:
\begin{enumerate}[leftmargin=2em]
\item Observed episode discharge $\mu_j$ against $z_j=D_j/C_j$ for the
Atlanta corridor.
\item Facility-specific constant-$\mu$ benchmark.
\item Fitted QVDF-indexed curve $\mu(z)$ with uncertainty band.
\item Observed congested vehicle speed compared with
$V_{\mathrm{FD}}^{\mathrm{cong}}[\mu(z)]$.
\item Out-of-sample comparison of fixed-$\mu$, episode-specific, and
QVDF-indexed emission estimates.
\end{enumerate}
\vspace{3mm}
\end{minipage}
}
\caption{Placeholder for independent validation of the QVDF-indexed
episode-equivalent discharge relation using the Atlanta corridor dataset.}
\label{fig:atlanta_validation_placeholder}
\end{figure}


% ============================================================
% BLOCK O
% INSERT IN THE CONCLUSION BEFORE:
% \paragraph{Limitations and scope.}
% ============================================================

The classical fixed bottleneck is nested in the proposed framework. With
constant discharge $\mu_0$, the FD-equivalent congested vehicle speed and
delay-to-emission factor are constant, and the incremental emission of a
vehicle is exactly proportional to its waiting time:
\begin{equation}
\Delta e(\tau)=\Gamma_0w(\tau).
\end{equation}
The QVDF extension retains this within-episode constancy but permits the
episode-equivalent discharge, congested vehicle speed, and emission severity
per unit delay to differ across representative congestion levels. Hence,
$D/C$ affects emissions through two separate channels: a delay-exposure
channel and a physical-state channel. This distinction is the central
theoretical contribution of the paper.

The present formulation should be interpreted as a reduced-form bridge from
period-aggregate assignment variables to bottleneck-consistent queue
emissions. It does not reconstruct a complete endogenous departure-time
equilibrium, queue-boundary propagation, moving bottlenecks, or the full
spatial heterogeneity of a long physical queue. Those extensions require a
dynamic traffic assignment or kinematic-wave formulation.
